\subsection{Reactive Type System of \GRSync}
\label{sec:semantics:grsync:react_typing}

While \GRSync components are synchronous, they are designed to be instantiated within the
asynchronous, event-driven services of \GRFRP. To ensure efficient execution while maintaining all
flows up-to-date, services in \GRFRP compute by propagating signal updates and event occurrences.
A key challenge in such a system is to ensure that service behavior remains \emph{compositional}:
extending a service with new imports or timers must not alter its original observable behavior.
Compositionality guarantees that outputs (exports and timers) are determined solely by declared
dependencies, not by the timing or arrival rate of inputs (imports and timers).

This guarantee is challenged when, for instance, a service is extended with a new import. A value
arriving on this new import triggers an execution of the service. We refer to this situation as an
\emph{oversampling} of the original service logic, since a computation occurs without any change to
the initial set of inputs. The risk is that this unintended execution could alter the internal state
of components or produce new outputs. Such side effects would change subsequent computations and
violate the principle of compositionality.

Three issues threaten this principle.

\paragraph{Three compositionality issues}

The first issue involves interaction with physical time. The \GRFRP language provides a built-in
operator, \timeop, that produces a flow representing the current time. The values of this flow are
sourced from the timestamps of all arriving inputs. This design creates a direct conflict with
compositionality.
%
Consider the \idf{integrate} component defined below:
\begin{equation*}
    \componentTab[\idf{x} = 0, \idf{t} = 0]{\idf{integrate}}{\idf{v}, \idf{t}}{\idf{x}}{
        \idf{x} = \last{\idf{x}} + \idf{v}*(\idf{t} - \last{\idf{t}})
    }
\end{equation*}
It integrates a flow \idf{v} over a flow \idf{t} representing the time. This component is intended
for use in a service where the input \idf{t} is connected to \timeop.
%
\Cref{tab:semantics:grsync:react_typing:last} demonstrates how oversampling corrupts the computation
by comparing a reference execution (left) with an oversampled one (right). The oversampling at
$\idf{t} = 6$ (\textcolor{gray}{gray}) changes the value of the \timeop operator, causing the
component recompute and alter its internal state. This additional computation causes all subsequent
outputs to diverge from the reference execution (in \textcolor{red}{red}).

\begin{table}[ht!]
    \centering
    \begin{tabular}{r|ccccc ccc r|cccccc}
        \idf{v}        & 2 & 3 & 3  & 2                    & 4                   &                     &  &  &
        \idf{v}        & 2 & 3 & 3  & \textcolor{gray}{3}  & 2                   & 4                           \\
        \idf{t}        & 1 & 3 & 4  & 7                    & 9                   &                     &  &  &
        \idf{t}        & 1 & 3 & 4  & \textcolor{gray}{6}  & 7                   & 9                           \\
        \last{\idf{t}} & 0 & 1 & 3  & 4                    & 7                   &                     &  &  &
        \last{\idf{t}} & 0 & 1 & 3  & \textcolor{gray}{4}  & \textcolor{red}{6}  & 7                           \\
        \idf{x}        & 2 & 8 & 11 & 17                   & 25                  &                     &  &  &
        \idf{x}        & 2 & 8 & 11 & \textcolor{gray}{17} & \textcolor{red}{19} & \textcolor{red}{27}         \\
        \last{\idf{x}} & 0 & 2 & 8  & 11                   & 17                  &                     &  &  &
        \last{\idf{x}} & 0 & 2 & 8  & \textcolor{gray}{11} & \textcolor{red}{17} & \textcolor{red}{19}         \\
    \end{tabular}
    \caption[A \kwf{last} memory breaks the oversampling.]%
    {Two executions with the same input flow: one reacting on input changes (left), and one
        oversampled at time $\idf{t} = 6$ in \textcolor{gray}{gray} (right). Oversampling caused
        modifications, in \textcolor{red}{red}, during the execution afterwards.}
    \label{tab:semantics:grsync:react_typing:last}
\end{table}

The second issue stems from the unrestricted use of \kwf{last} memories. Consider the following
\idf{count} component, which implements a simple counter with no inputs:
\begin{equation*}
    \component{\idf{count}}{}{\idf{o}}{\idf{o} = \last{\idf{o}} + 1}{\idf{o} = 0}.
\end{equation*}
From an event-driven perspective, this component should remain inactive, as it has no inputs to
trigger a computation. However, its definition implies a free-running feedback loop through the
memory \last{\idf{o}}: the value of \idf{o} at any given instant is used to define the value of
\last{\idf{o}} at the very next instant. The output flow is therefore \emph{continuous} in the sense
that its definition suggests it should be continuously re-evaluated in order to maintain its value
up-to-date. This ambiguity makes the behavior dependent on an implicit execution rate rather than a
reaction to explicit inputs, violating the principles of reactive programming.

The third issue arises from unrestricted event emissions, using the \kwf{emit} operator. Consider
the \idf{emit\_now} component, which creates an event that unconditionally occurs:
\begin{equation*}
    \component{\idf{emit\_now}}{\idf{s}}{\idf{e}}{\idf{e} = \emit{\idf{s}}}{}.
\end{equation*}
Conceptually, the output event \idf{e} is \emph{continuously} present, since its emission is not
guarded by any condition. An event-driven execution, however, would only trigger this component when
the input signal \idf{s} is updated, emitting the event only at that time. This creates a
fundamental conflict between the declarative semantics and the reactive execution model.

\paragraph{The solution: explicit activation}

To resolve these issues, \GRust provides the \kwf{when} construct, which enables computations and
event emissions to be triggered explicitly. It activates a computation block only upon the
occurrence of a discrete event, such as the occurrence of an event or the rising edge of a boolean
signal. This mechanism ensures that computations are executed only at precise, explicitly defined
moments.

The first issue, caused by the implicit dependency on \timeop, is resolved by guarding the
integration. The \idf{integrate\_every} component updates the position \idf{x} only when the event
\idf{i} occurs:
\begin{equation*}
    \componentTab{\idf{integrate\_every}}{\idf{v}, \idf{t}, \idf{i}}{\idf{x}}{
        \whenTab{\idf{x}}{\idf{x} = 0}{
            \branch{\detect{\_}{\idf{i}}}{\true}{\idf{x} = \idf{integrate}(\idf{v}, \idf{t})}
        }
    }
\end{equation*}
Because the \kwf{when} clause conditions the execution, the output \idf{x} is not updated during
oversampling, \emph{provided the event \idf{i} does not occur during oversampling} (we come back to
this point in Paragraph~\textbf{iii}). This restores compositionality.

Similarly, the \kwf{when} construct addresses the ambiguity of the free-running counter. The
\idf{count\_every} component increments its output only upon the occurrence of event \idf{i}:
\begin{equation*}
    \componentTab{\idf{count\_every}}{\idf{i}}{\idf{o}}{
        \whenTab{\idf{o}}{\idf{o} = 0}{
            \branch{\detect{\_}{\idf{i}}}{\true}{\idf{o} = \idf{count}()}
        }
    }
\end{equation*}
This approach discretizes the behavior of the \idf{count} component, giving it a clear, reactive
semantic. The \idf{count\_every} component is therefore resilient to oversampling \emph{as long as
    the triggering event \idf{i} is absent during oversampling}.

Finally, the problem of continuous event emission is handled in the same manner. The component
\idf{emit\_every} emits an event only when the event \idf{i} occurs:
\begin{equation*}
    \componentTab{\idf{emit\_every}}{\idf{s}, \idf{i}}{\idf{e}}{
        \whenTab{\idf{e}}{}{
            \branch{\detect{\_}{\idf{i}}}{\true}{\idf{e} = \idf{emit\_now}(\idf{s})}
        }
    }
\end{equation*}
Consequently, the event \idf{e} is emitted only at explicitly defined moments, preventing unintended
emissions during oversampling, \emph{as long as the event \idf{i} is absent during oversampling},
and aligning the behavior with the event-driven execution model.

\paragraph{A reactive type system}

In the preceding examples, the correctness of the \kwf{when} construct hinges on a critical
assumption: \emph{the triggering event \idf{i} must not occur during an oversampling execution}.
\GRust enforces this property statically through a reactive type system. This system analyzes the
activation conditions of every flow to ensure it is only activated by its declared dependencies. By
doing so, it restricts the influence of the three sources of compositionality issues: \timeop,
unguarded \kwf{last} memories, and unconditional event emissions. As a result, oversampling by an
unrelated input cannot produce spurious events or rising edges, preserving compositionality by
construction.

To formalize this guarantee, the type system distinguishes between two kinds of flows, based on
their sensitivity to oversampling:
%
\begin{itemize}
    \item \emph{Continuous flows}, which are inherently tied to the execution step itself. This
          category includes \timeop, as well as unguarded \kwf{last} memories and event emissions.
          Their values can change with every computation, making them sensitive to the execution
          rate.

    \item \emph{Discrete flows}, which encompass all other flows. These are driven by explicit input
          events and are guaranteed to be stable during oversampling: signals hold their last value,
          and events remain absent.
\end{itemize}

The reactive type system classifies every flow as either continuous or discrete, statically
prohibiting the activation of discrete computations by continuous sources. This section details the
\GRSync portion of this system. By enforcing these rules, the type system provides the foundation
for the oversampling theorem (\Cref{thm:semantics:grsync:theorems:oversampling}), which guarantees
that event-driven execution preserves the semantics of the underlying synchronous model.

\begin{table}[!ht]
    \centering
    \begin{tabularx}{\textwidth}{|l|X|}\hline
        \rType ::= \cType $\vert$ \dType                       & Reactive type denoting if the
        associated element is continuous (depending on \timeop or untriggered event emissions and
        \kwf{last} memories) or discrete.                                                            \\\hline
        $\Theta[x] = \rType$                                   & A reactive typing context $\Theta$
        mapping flow identifiers to their \rType.                                                    \\\hline
        \compTypedReact{f}{\rType^+}{\rType^+}                 & The component $f$ is reactive and,
        when given input types $\rType^+$, it returns values of types $\rType^+$.                    \\\hline
        \eqTypedReact{\Theta}{eq^+}                            & The equations $eq^+$ are reactive
        in the reactive typing context $\Theta$.                                                     \\\hline
        \exprTypedReact{\Theta}{e}{\rType^+}                   & The expression $e$ has type
        $\rType^+$ in the reactive typing context $\Theta$.                                          \\\hline
        \patTypedReact{\Theta}{\guarded{pat^+}{e_g}}{\rType^+} & The pattern \guarded{pat^+}{e_g} is
        reactive to the type $\rType^+$ in the reactive typing context $\Theta$.                     \\\hline
        \epatTypedReact{\Theta}{\guarded{epat^+}{e_g}}         & The guarded event pattern
        \guarded{epat^+}{e_g} is reactive in the reactive typing context $\Theta$.                   \\\hline
    \end{tabularx}
    \caption[Notations for the reactive type system of \GRSync.]%
    {Notations for the logic of the reactive type system of \GRSync.}
    \label{tab:semantics:grsync:react_typing:notations}
\end{table}

We first introduce the notations used throughout this type system, summarized in
\Cref{tab:semantics:grsync:react_typing:notations}.
%
The reactive type \rType is used to detect a dependency to the \timeop operation and to an
untriggered event emission and \kwf{last} memory. It is recorded in a reactive typing context
$\Theta$, where for all identifier $x$, $\Theta[x] = \cType$ denotes that $x$ depends on those
constructs ($x$ is a \emph{continuous} flow), and $\Theta[x] = \dType$ denotes that $x$ has no such
dependencies ($x$ is a \emph{discrete} flow).
%
When given the input types $\rType_i^+$, the component $f$ is reactive and returns values of type
$\rType_o^+$ if we have \compTypedReact{f}{\rType_i^+}{\rType_o^+}.
%
The judgement \eqTypedReact{\Theta}{eq^+} denotes that the equations $eq^+$ are reactive and that
the typing context $\Theta$ corresponds to the induced dependencies.
%
An expression $e$ has type $\rType^+$ in the typing context $\Theta$ if we have
\exprTypedReact{\Theta}{e}{\rType^+}.
%
When matching an expression of type $\rType^+$, the pattern \guarded{pat^+}{e_g} is reactive in
the reactive typing context $\Theta$ if we have \patTypedReact{\Theta}{\guarded{pat^+}{e_g}}
{\rType^+}.
%
The judgement \epatTypedReact{\Theta}{\guarded{epat^+}{e_g}} denotes that the event pattern
\guarded{epat^+}{e_g} is reactive and that the typing context $\Theta$ corresponds to the induced
dependencies.

The following paragraphs explain in details the rules that define the \GRSync portion of this
reactive type system.

\paragraph{Components}

The reactivity of a component is crucial to ensure a sound usage in an asynchronous service.
\Cref{thm:semantics:grsync:theorems:oversampling} states that a reactive component can be executed
in an event-driven manner, only propagating changes. This execution is proved to be resilient to
oversampling: the absence of changes from the discrete input flows of the component leads to an
absence of changes from its internal memory and from its discrete output flows. The presented
reactive type system tracks dependencies to continuous constructs inside \emph{top-level
    components} (\ie components that are called by services). Components that are not reactive must
be called in reactive components---that are in turn called by services. The tracking ends as soon as
it enters a \kwf{when} equation. This kind of equation discretizes computations, preventing changes
during an oversampling.

Given the input and output reactive types $\rType_i^+$ and $\rType_o^+$, we say that the component $f$
(\component{f}{x_i^+}{x_o^+}{eq^+}{x_m^+ = c^+}) is reactive, following the rule \nameref{rule:%
    semantics:grsync:react_typing:comp}, if there exists a typing context $\Theta$ such that it maps
the input and output identifiers to their corresponding types ($\Theta[x_i] = \rType_i$ and
$\Theta[x_o] = \rType_o$), the memorized identifiers are discrete ($\Theta[x_m] = \dType$), their
memory identifiers are continuous ($\Theta[\last{x_m}] = \cType$), and the block of equations
is reactive. The memorized identifiers need to be discrete to preserve the state of the component
during an oversampling: if they were not, their values could change during an oversampling and be
stored in memory for the next computations.

\begin{center}
    \begin{minipage}{0.95\textwidth}
        \begin{mathpar}
            \inferREACT
            {
                G[f] = \component{f}{x_i^+}{x_o^+}{eq^+}{x_m^+ = c^+} \\
                \eqTypedReact{\Theta}{eq^+} \\
                \left(\Theta[x_i] = \rType_i\right)^+ \\
                \left(\Theta[x_o] = \rType_o\right)^+ \\
                \left(\Theta[x_m] = \dType \right)^+ \\
                \left(\Theta[\last{x_m}] = \cType \right)^+
            }
            { \compTypedReact{f}{\rType_i^+}{\rType_o^+} }
            {Comp}
            \label{rule:semantics:grsync:react_typing:comp}
        \end{mathpar}
    \end{minipage}
\end{center}

This rule allows components to return continuous flows. For instance, component that multiplies
\timeop by two is not an issue as long as its output is tracked as continuous.

\paragraph{Expressions}

The reactive type system for expressions informs whether the expression can change during an
oversampling (\ie the expression is continuously), or if it is resilient to it (\ie the expression
is discrete). It propagates the continuous dependencies.
%
The two continuous expressions (\last{x_m} and \emit{e}) are propagated by the
rules \nameref{rule:semantics:grsync:react_typing:e:last} and
\nameref{rule:semantics:grsync:react_typing:e:emit}. For the \kwf{last} memory, the type stored in
the context is verified to be the continuous type~\cType.

\begin{center}
    \begin{minipage}{0.9\textwidth}
        \begin{mathpar}
            \inferREACTE
            { \Theta[\last{x_m}] = \cType }
            { \exprTypedReact{\Theta}{\last{x_m}}{\cType} }
            {Last}
            \label{rule:semantics:grsync:react_typing:e:last}

            \inferREACTE
            { }
            { \exprTypedReact{\Theta}{\emit{e}}{\cType} }
            {Emit}
            \label{rule:semantics:grsync:react_typing:e:emit}
        \end{mathpar}
    \end{minipage}
\end{center}

% For example, let $f$ be a component with input \idf{x_i} and output \idf{x_o}, such that \init
% {\idf{x_o}= 0} and $\idf{x_o} = \last{\idf{x_o}} + \idf{x_i}$. A call of $f$ in a service would
% define a flow whose values depend on the execution times and not only on value updates. This is
% depicted in the chronograms from \Cref{tab:semantics:grsync:react_typing:last}, illustrating two
% executions starting with the same input flow: one execution reacting on the input changes (left),
% and the other oversampled at time $t_{\text{OS}}$ (right). Oversampling caused a difference of
% output value for the timestamp $t_2$.

% \begin{table}[ht!]
%     \centering
%     \begin{tabular}{r|llll ccc r|lllll}
%         Time                                       & $t_0$ & $t_1$ & $t_2$                             & \dots &       &  &  &
%         Time                                       & $t_0$ & $t_1$ & \textcolor{gray}{$t_{\text{OS}}$} & $t_2$ & \dots         \\
%         \idf{x_i}                                  & $0$   & $1$   & $2$                               & \dots &       &  &  &
%         \idf{x_i}                                  & $0$   & $1$   & \textcolor{gray}{$1$}             & $2$   & \dots         \\
%         $\idf{x_o} = \last{\idf{x_o}} + \idf{x_i}$ & $0$   & $1$   & $3$                               & \dots &       &  &  &
%         $\idf{x_o} = \last{\idf{x_o}} + \idf{x_i}$ & $0$   & $1$   & \textcolor{gray}{$2$}             & $4$   & \dots         \\
%     \end{tabular}
%     \caption{Two executions with the same input flow: one reacting on input changes (left), and one
%         oversampled at time $t_{\text{OS}}$ (right). Oversampling caused a shift for the timestamp
%         $t_2$.}
%     \label{tab:semantics:grsync:react_typing:last}
% \end{table}

% Similarly, an event emission outside the branch of a \kwf{when} equation could simulate an event
% arrival, allowing the usage of \kwf{last}. For example, the following equations present the same
% issue depicted by the chronograms from \Cref{tab:semantics:grsync:react_typing:last}.
% \begin{align*}
%      & \idf{x} = \emit{()},                                                                            \\
%      & \when{\idf{y}}{\idf{y} = 0}{\branch{\detect{\_}{\idf{x}}}{\true}{\idf{y} = \last{\idf{y}} + 1}}
% \end{align*}

The rule \nameref{rule:semantics:grsync:react_typing:e:const} ensures constants ($c$) are of type
\dType, as they do not depend on any continuous flow.
%
Identifiers ($x$) have the type stored in the context $\Theta$, as stated by \nameref{rule%
    :semantics:grsync:react_typing:e:ident}.

\begin{center}
    \begin{minipage}{0.9\textwidth}
        \begin{mathpar}
            \inferREACTE
            { }
            { \exprTypedReact{\Theta}{c}{\dType} }
            {Cst}
            \label{rule:semantics:grsync:react_typing:e:const}

            \inferREACTE
            { }
            { \exprTypedReact{\Theta}{x}{\Theta[x]} }
            {ID}
            \label{rule:semantics:grsync:react_typing:e:ident}
        \end{mathpar}
    \end{minipage}
\end{center}

An operation ($op(e^+)$) is of type \dType if all input expressions ($e^+$) are of type
\dType, as defined by the rule \nameref{rule:semantics:grsync:react_typing:e:op:value}.
In the case one expression has type \cType, the rule \nameref{rule:semantics:grsync:react_typing:%
    e:op:time} states the operation is of type \cType. Those rules implement the propagation of
continuous dependencies.

\begin{center}
    \begin{minipage}{0.9\textwidth}
        \begin{mathpar}
            \inferREACTE
            { \left(\exprTypedReact{\Theta}{e}{\dType}\right)^+ }
            { \exprTypedReact{\Theta}{op(e^+)}{\dType} }
            {Op{\scriptsize Disc}}
            \label{rule:semantics:grsync:react_typing:e:op:value}

            \inferREACTE
            {
                \left(\exprTypedReact{\Theta}{e}{\rType}\right)^+ \\\\
                \exists \rType = \cType
            }
            { \exprTypedReact{\Theta}{op(e^+)}{\cType} }
            {Op{\scriptsize Cont}}
            \label{rule:semantics:grsync:react_typing:e:op:time}
        \end{mathpar}
    \end{minipage}
\end{center}

A component application $f(e^+)$ is typed $\rType_o^+$ if the input expressions are typed
$\rType_i^+$, and if the called component ($f$) is reactive with the input types $\rType_i^+$ and
the output types $\rType_o^+$. This is defined by the rule
\nameref{rule:semantics:grsync:react_typing:e:app}. As previously said, the reactive type system
tracks and propagates dependencies to continuous constructs inside top-level components. This
propagation ends as soon as it enters a \kwf{when} equation---which is resilient to oversampling by
construction. Therefore, the reactive typing of expressions ends in \kwf{when} equations; the rule
\nameref{rule:semantics:grsync:react_typing:e:app} is only applied outside \kwf{when} equations, and
thus requires the called component ($f$) to be reactive.

\begin{center}
    \begin{minipage}{0.9\textwidth}
        \begin{mathpar}
            \inferREACTE
            {
                \left(\exprTypedReact{\Theta}{e}{\rType_i}\right)^+ \\
                \compTypedReact{f}{\rType_i^+}{\rType_o^+}
            }
            { \exprTypedReact{\Theta}{f(e^+)}{\rType_o^+} }
            {App}
            \label{rule:semantics:grsync:react_typing:e:app}
        \end{mathpar}
    \end{minipage}
\end{center}

\paragraph{Patterns}

A pattern matching is reactive if and only if an oversampling does not change the matched branch. So
it is reactive if and only if the success of a pattern is resilient to oversampling. The reactivity
of a pattern captures this resilience.
%
Guarded patterns (\guarded{pat^+}{e_g}) are reactive, with regard to the matching types
($\rType^+$), if each individual pattern ($pat$) is reactive with its corresponding type, and if the
guard ($e_g$) is of discrete type \dType, as stated by the rule
\nameref{rule:semantics:grsync:react_typing:pats}. The guard should be discrete so that its value
does not change during an oversampling.

\begin{center}
    \begin{minipage}{0.9\textwidth}
        \begin{mathpar}
            \inferREACTPat
            {
                \left(\patTypedReact{\Theta}{pat}{\rType}\right)^+ \\
                \exprTypedReact{\Theta}{e_g}{\dType}
            }
            { \patTypedReact{\Theta}{\guarded{pat^+}{e_g}}{\rType^+} }
            {Tple}
            \label{rule:semantics:grsync:react_typing:pats}
        \end{mathpar}
    \end{minipage}
\end{center}

The rule \nameref{rule:semantics:grsync:react_typing:pat:const} ensures that a constant pattern is
reactive if and only if it matches a discrete type \dType. Indeed, a constant pattern acts just like
a guard, so we need to make sure that the matched expression does not change its value during the
oversampling.

\begin{center}
    \begin{minipage}{0.9\textwidth}
        \begin{mathpar}
            \inferREACTPat
            { }
            { \patTypedReact{\Theta}{c}{\dType} }
            {Const}
            \label{rule:semantics:grsync:react_typing:pat:const}
        \end{mathpar}
    \end{minipage}
\end{center}

Without this restriction, all identifiers defined by a \kwf{match} equation could depend on \timeop.
For example, if \idf{t} is an identifier depending on \timeop, then the following block of equations
presents the issue depicted by the chronograms from \Cref{tab:semantics:grsync:react_typing:pat}:
the selected branch in the pattern matching changes during the oversampling, making the rising edge
on \idf{x} tick, which in turn increments \idf{f}, altering the memory of the program for the long
term.
\begin{equation*}
    \left\{
    \matchTab{\idf{t}}{\idf{x}}{\branch{6}{\true}{\idf{x} = \true},\\\qquad\branch{\_}{\true}{\false}}
    \qquad
    \whenTab{\idf{y}}{\idf{y} = 0}{\branch{\idf{x}}{\true}{\idf{y} = \last{\idf{y}} + 1}}
    \right\}
\end{equation*}
%
The reactive type system prevents this issue: \idf{t} is associated with a \cType by rule
\nameref{rule:semantics:grsync:react_typing:e:ident}, the \kwf{match} equation follows rule
\nameref{rule:semantics:grsync:react_typing:eq:match} and tries to associate the constant pattern
\guarded{6}{\true} to the \cType type, but only \dType type is allowed for this pattern, as stated
by rule \nameref{rule:semantics:grsync:react_typing:pat:const}.

\begin{table}[ht!]
    \centering
    \begin{tabular}{r|ccc ccc r|cccc}
        \idf{t}        & 1      & 7                       & 9                  &                    &  &  &
        \idf{t}        & 1      & \textcolor{gray}{6}     & 7                  & 9                          \\
        \idf{x}        & \false & \false                  & \false             &                    &  &  &
        \idf{x}        & \false & \textcolor{gray}{\true} & \false             & \false                     \\
        \idf{y}        & 0      & 0                       & 0                  &                    &  &  &
        \idf{y}        & 0      & \textcolor{gray}{1}     & \textcolor{red}{1} & \textcolor{red}{1}         \\
        \last{\idf{y}} & 0      & 0                       & 0                  &                    &  &  &
        \last{\idf{y}} & 0      & \textcolor{gray}{0}     & \textcolor{red}{1} & \textcolor{red}{1}         \\
    \end{tabular}
    \caption[A \kwf{match} equation breaks the oversampling.]%
    {Two executions with the same input flow: the left one is the reference execution, and
        the right one is oversampled at time $\idf{t} = 6$ (in \textcolor{gray}{gray}). Oversampling
        caused modifications, in \textcolor{red}{red}, during the execution afterwards.}
    \label{tab:semantics:grsync:react_typing:pat}
\end{table}

An identifier pattern ($x$), following the rule \nameref{rule:semantics:grsync:react_typing:pat:%
    ident}, associates to the identifier the corresponding type stored in the context ($\Theta[x]$).
%
A wildcard is reactive and corresponds to any type, as stated by the rule \nameref{rule:semantics:%
    grsync:react_typing:pat:wild}, because it always matches (with or without oversampling).

\begin{center}
    \begin{minipage}{0.9\textwidth}
        \begin{mathpar}
            \inferREACTPat
            { }
            { \patTypedReact{\Theta}{x}{\Theta[x]} }
            {ID}
            \label{rule:semantics:grsync:react_typing:pat:ident}

            \inferREACTPat
            { }
            { \patTypedReact{\Theta}{\_}{\rType} }
            {Wild}
            \label{rule:semantics:grsync:react_typing:pat:wild}
        \end{mathpar}
    \end{minipage}
\end{center}

\paragraph{Event patterns}

Similarly to pattern matching, a \kwf{when} equation is reactive if and only if an oversampling does
not change the triggered branch. So it is reactive if and only if the success of an event pattern is
resilient to oversampling.
%
Guarded event patterns (\guarded{epat^+}{e_g}) are reactive if each individual event pattern
($epat$) is reactive, as stated by the rule \nameref{rule:semantics:grsync:react_typing:epats}.
Because each individual event pattern is resilient to oversampling (is not triggered during an
oversampling), the guard is computed only when a change occurred, and thus can depend on continuous
flows.

\begin{center}
    \begin{minipage}{0.9\textwidth}
        \begin{mathpar}
            \inferREACTEPat
            { \left(\epatTypedReact{\Theta}{epat}\right)^+ }
            { \epatTypedReact{\Theta}{\guarded{epat^+}{e_g}} }
            {Tple}
            \label{rule:semantics:grsync:react_typing:epats}
        \end{mathpar}
    \end{minipage}
\end{center}

Event detections are reactive if the event (and the identifiers they define) are of discrete type
\dType, as stated by the rule \nameref{rule:semantics:grsync:react_typing:epat:event}.
%
The detection of a rising edge ($e$) follows the rule \nameref{rule:semantics:grsync:react_typing:%
    epat:edge}, and is reactive if and only if the boolean expression ($e$) is of type \dType.
This makes sure that only discrete changes are detected by an event pattern.

\begin{center}
    \begin{minipage}{0.9\textwidth}
        \begin{mathpar}
            \inferREACTEPat
            { \Theta[y] = \dType \\ \Theta[x] = \dType }
            { \epatTypedReact{\Theta}{\detect{x}{y}} }
            {Dtct}
            \label{rule:semantics:grsync:react_typing:epat:event}

            \inferREACTEPat
            { \exprTypedReact{\Theta}{e}{\dType} }
            { \epatTypedReact{\Theta}{e} }
            {Edge}
            \label{rule:semantics:grsync:react_typing:epat:edge}
        \end{mathpar}
    \end{minipage}
\end{center}

\newpage

\paragraph{Equations}
%
Equations link expressions dependencies to identifiers. They are reactive with regard to the
reactive type context $\Theta$ if the dependencies to continuous flows that they induce satisfy the
typing context.
%
A block of equations ($eq^+$) is reactive if and only if each individual equation ($eq$) is
reactive, as defined by the rule \nameref{rule:semantics:grsync:react_typing:eqs}.
%
A declaration ($x^+ = e$) is reactive and satisfies the typing context ($\Theta$), following the
rule \nameref{rule:semantics:grsync:react_typing:eq:decl}, if the reactive type of the expression
($e$) corresponds to the context ($\Theta[x]$).

\begin{center}
    \begin{minipage}{0.9\textwidth}
        \begin{mathpar}
            \inferREACTEq
            { \left(\eqTypedReact{\Theta}{eq}\right)^+ }
            { \eqTypedReact{\Theta}{eq^+} }
            {Blck}
            \label{rule:semantics:grsync:react_typing:eqs}

            \inferREACTEq
            { \exprTypedReact{\Theta}{e}{(\Theta[x])^+} }
            { \eqTypedReact{\Theta}{x^+ = e} }
            {Decl}
            \label{rule:semantics:grsync:react_typing:eq:decl}
        \end{mathpar}
    \end{minipage}
\end{center}

A \kwf{match} equation (\match{e}{x^+}{(\branch{pat_k^+}{e_k}{eq_k^+})^+}) is reactive if the
reactive types of the matched expression ($e$) make the patterns (\guarded{pat_k^+}{e_k}) reactive
in the context $\Theta$, and each branch equations ($eq_k^+$) are reactive, as defined by the rule
\nameref{rule:semantics:grsync:react_typing:eq:match}.

\begin{center}
    \begin{minipage}{0.95\textwidth}
        \begin{mathpar}
            \inferREACTEq
            {
                \exprTypedReact{\Theta}{e}{\rType^+} \\
                \left(\patTypedReact{\Theta}{\guarded{pat_k^+}{e_k}}{\rType^+}\right)^+ \\
                \left(\eqTypedReact{\Theta}{eq_k^+}\right)^+
            }
            { \eqTypedReact{\Theta}{\match{e}{x^+}{(\branch{pat_k^+}{e_k}{eq_k^+})^+}} }
            {Mtch}
            \label{rule:semantics:grsync:react_typing:eq:match}
        \end{mathpar}
    \end{minipage}
\end{center}

A \kwf{when} equation \emph{discretizes continuous flows} in reactive components. The rule
\nameref{rule:semantics:grsync:react_typing:eq:when} ensures \when{x^+}{x_m^+ = c^+}{(\branch
    {epat_k^+}{e_k}{eq_k^+})^+} is reactive if and only if its event patterns
(\guarded{epat_k^+}{e_k}) are reactive, meaning that the event detection is resilient to
oversampling. All defined identifiers ($x^+$) are of discrete type \dType. The reactive type system
ends as soon as it enters the \kwf{when} equation, because everything inside is discrete.

\begin{center}
    \begin{minipage}{0.9\textwidth}
        \begin{mathpar}
            \inferREACTEq
            {
                \left(\Theta[x] = \dType\right)^+ \\
                \left(\epatTypedReact{\Theta}{\guarded{epat_k^+}{e_k}}\right)^+
            }
            { \eqTypedReact{\Theta}{\when{x^+}{x_m^+ = c^+}{(\branch{epat_k^+}{e_k}{eq_k^+})^+}} }
            {When}
            \label{rule:semantics:grsync:react_typing:eq:when}
        \end{mathpar}
    \end{minipage}
\end{center}

\paragraph{Relation with the type system of \Zelus}

The reactive type system introduced in this section is reminiscent of the approach taken by the
synchronous language \Zelus~\cite{DBLP:conf/emsoft/BenvenisteBCP11} to solve a similar problem. In \Zelus, a language
to model hybrid systems, the core challenge is managing the interaction between discrete software
logic and continuous physical dynamics described by Ordinary Differential Equations (ODEs). A
synchronous program operates in discrete steps, with its state defined only at these specific
instants. In contrast, the values described by ODEs evolve continuously, changing smoothly between
these steps. This creates a fundamental conflict: the discrete program only sees samples of the
continuous world, making it blind to crucial events (like a value crossing a threshold) that could
occur between samples. Consequently, operations like checking for equality on a continuous value are
ill-defined in the discrete logic.

To ensure predictable behavior, the \Zelus type system enforces a strict separation between its
\emph{continuous values} (derived from ODEs) and \emph{discrete values}. This guarantees the
soundness of their interaction by preventing such ill-defined operations and ensuring the continuous
domain can only influence the discrete logic at specific, well-defined moments (such as
zero-crossing events).

While \GRust does not model ODEs, its \timeop operator introduces a similar conflict. It provides a
continuous flow that changes at every evaluation, clashing with the discrete flows during an
oversampling step. Our reactive type system classifies flows as either continuous (typed \cType) or
discrete (typed \dType) to ensure that continuous time from \timeop cannot improperly influence stateful
computations, thereby ensuring the soundness of the program during an oversampling.

\paragraph{Relation with Call-By-Push-Value}

A more theoretical understanding of our reactive type system is provided by Call-By-Push-Value
(CBPV)~\cite{DBLP:phd/ethos/Levy01}. CBPV is a programming paradigm that distinguishes \emph{value
    types} from \emph{computation types}. Value types represent stable data that can be passed and
stored. In contrast, computation types represent actions that, when executed, produce a value and
may have effects.

In CBPV, to treat a computation as a value, it must be packaged into a \emph{thunk} that freezes
the computation into a value. To run the computation, the thunk must be explicitly \emph{forced}.
This formalizes the distinction between having data \versus having a process to generate data.

Our reactive type system can be understood as a domain-specific application of this principle. A
discrete flow (typed \dType) corresponds to a CBPV value: it represents stable data that is
unaffected by re-evaluation during an oversampling step. Conversely, a continuous flow
(typed \cType), originating from \timeop, corresponds to a CBPV computation: evaluating \timeop is
an action (getting the current time) that produces a new result each time.

From this perspective, our reactive type system acts as a static guard that restricts where the
implicit \emph{force} operation on \timeop is allowed to occur. It forbids forcing a time-based
computation if its result could alter the discrete memory of the program.

\paragraph{Summary}

This section has introduced the reactive type system for \GRSync. The primary goal of this system is
to resolve the conflict between the synchronous semantics of components and the asynchronous,
event-driven execution model of services. It formally addresses this challenge by distinguishing
between \emph{continuous} flows, which depend on execution time or implicit feedback loops, and
\emph{discrete} flows, which are driven by explicit input changes.

The presented rules systematically propagate continuity, ensuring that dependencies on \timeop and
untriggered \kwf{last} memories and \kwf{emit} operations are tracked. The \kwf{when} equation is
the central mechanism for discretizing this behavior. By enforcing that its triggers are discrete,
the \kwf{when} construct establishes a boundary that prevents continuous flows from causing
unintended effects during oversampling.

A component that adheres to this type system is consequently deemed \emph{reactive}. This static
guarantee is essential, as it ensures the component can be safely and compositionally integrated
into an asynchronous system (\cf \Cref{thm:semantics:grsync:theorems:oversampling}), providing a
sound foundation for the event-driven execution model of \GRust.
